% \chapter{Axionlike-particle-induced spin precession}\label{chap:theory}

\section{Thermalized Milky Way ALP}\label{sec:theory:MWALP}

\YZ{ideal gas treatment: discuss thermalized halo and velocity distribution.}

% estimation of the Rabi frequency $\Omega_a^\mathrm{Earth}$ for the Earth-bound ALP

The r.m.s.\ Rabi frequency $\Omega_a$ can be computed based on Equations~\eqref{eq:ALP_field}, \eqref{eq:a0_rhoDM}, and~\eqref{eq:HaNN} \cite{Centers2021Stochastic}:
\begin{equation}
    \Omega_a = \dfrac{1}{2\hbar} \mathrm{g_{aNN}} a_0 m_a c v_{\bot} \approx \dfrac{1}{2} \mathrm{g_{aNN}} \sqrt{2\hbar c \rho_{DM}} v_{\bot}
    \,,
    \label{eq:Omega_a_MW}
\end{equation}
%\AW{this is a stochastically varying quantity!}
Here we assume that the ALP-field gradient comes exclusively from the relative motion of the field and the detector.\footnote{This is not necessarily the case, for instance, for ALPs gravitationally bound in stationary states like axion stars or ``axion atoms'' \cite{Lam2017_PhysRevD, Liang_Zhitnitsky2019_GravitationallyPhysRevD}. We do not consider such regimes here.}
%\DK{The readers probably don't know what you're referencing here... perhaps use the term that appears more commonly in the literature, "axion stars," and include related references, for example:\cite{Braaten_2019_axionstars_RevModPhys, Eby_2019_QCDaxionstars_PhysRevD, Eby_2021_axionstarsPlanck_PhysRevD, Banerjee2020_Relaxionstars,Banerjee2020axionhalos}}
%\DB{But, indeed, we were not thinking about axion stars that are bound states of ALPs themselves, but rather really Axion Atom with, say, the Sun in the center and a halo around it. The atomic analogy is quite strong here and this has been used in the literature…}
% It is important to note that the ALP field is a stochastically varying quantity, so that the $\Omega_a$ in Equation~\eqref{eq:Omega_a} is the r.m.s value.
If we take $v_{\bot} = \SI{220}{\km\per\second}$, we have $\Omega_a/2\pi=\mathrm{g_{aNN}}\times\SI{0.2}{\GeV\cdot\Hz}$. For an experiment searching for $\mathrm{g_{aNN}}\leq 10^{-10}\,\SI{}{\per\GeV}$, the Rabi period is $2\pi/\Omega_a \geq 1000\,\mathrm{yr}$; the ALP-field gradient is therefore a weak drive on the spins.


\subsection{Coherence time}\label{sec:theory:MWALP:coherence}  % for MW axion

For a Maxwell-Boltzmann velocity distribution with average speed $v_a \approx \SI{220}{\km\per\s}$ \cite{Evans2019_SHMRefined}, the ALP field spectral linewidth (FWHM) observed by a detector on Earth is
\begin{equation}
    \Gamma_a = \dfrac{\nu_a}{Q_a}
    \,,\label{eq:Gamma_a_Qa}
\end{equation}
where $Q_a \equiv (c/v_a)^2 \approx 1.8\times10^6$ is the ALP quality factor. The ALP coherence time $\tau_a$, defined by $\Gamma_a = 2\pi/\tau_a$, is therefore
\begin{equation}
    \tau_a = \dfrac{2\pi Q_a}{\nu_a}
    \,.
    \label{eq:taua}
\end{equation}
Within one coherence time the superposition Equation~\eqref{eq:ALP_sum} is well approximated by the single coherent mode Equation~\eqref{eq:ALP_field}. On longer time scales the random phases between modes cause the amplitude and phase of the effective oscillation to drift, so the field must be treated as a stochastic process \cite{Centers2021Stochastic}. For Compton frequencies above \SI{1}{\kHz}, $\tau_a \lesssim \SI{1000}{\second}$; the interplay between $\tau_a$ and the NMR relaxation times determines the optimal measurement strategy in Chapter~\ref{chap:results}.




\subsection{Periodic modulations}\label{sec:theory:MWALP:modulations}

% \subsection{Daily modulation}

The ALP-field gradient driving spin precession depends on the direction of the relative velocity between the Earth and the local dark-matter halo. As the Earth rotates, this direction sweeps through space with a period of one sidereal day. An ALP signal therefore exhibits a daily modulation of its amplitude and phase \cite{Gramolin2022Feb_SpectralSignatures}.

% \subsection{Annual modulation}

In addition to daily rotation, the Earth's orbital velocity around the Sun modulates the relative ALP-detector velocity with a period of one year. Both the daily and annual modulations have precisely predictable periods and phases, making them useful consistency checks for any claimed dark-matter-candidate signal \cite{Gramolin2022Feb_SpectralSignatures}.




\subsection{Numerical simulation}\label{sec:theory:MW:simulation}

\YZ{include several important examples and give a summary of signal dependence. }
